\begin{exercise}[Matter--radiation equality]{Tong Cosmology, Cambridge Part III Sheet 3}
Using present photon, neutrino, and matter densities, derive the redshift and
comoving horizon scale at matter--radiation equality.  Explain why this scale
marks the turnover of the matter power spectrum.
\end{exercise}

\begin{exercise}[Entropy conservation and changing degrees of freedom]{Tong Cosmology, Cambridge Part III Sheet 3}
Show that adiabatic expansion conserves \(s a^3\), with
\(s=(2\pi^2/45)g_{\ast s}T^3\).  Derive the temperature scaling when
\(g_{\ast s}\) changes.
\end{exercise}

\begin{exercise}[The cosmic neutrino temperature]{Cambridge Part III Cosmology, 2015 exam}
Assuming instantaneous neutrino decoupling before electron--positron
annihilation, use entropy conservation to derive \(T_\nu/T_\gamma\) after
annihilation and the standard neutrino contribution to the radiation density.
\end{exercise}

\begin{exercise}[Baryon-to-photon ratio]{Tong Cosmology, Cambridge Part III Sheet 3}
Relate \(\eta=n_b/n_\gamma\) to \(\Omega_bh^2\) and the measured CMB
temperature.  Explain why \(\eta\) remains constant after entropy production
ceases and why it is central to nucleosynthesis.
\end{exercise}

\begin{exercise}[The sound horizon at recombination]{Cambridge Part III Cosmology, examples}
For the tightly coupled photon--baryon fluid, derive
\[
c_s^2=\frac{1}{3(1+R)},\qquad R=\frac{3\rho_b}{4\rho_\gamma},
\]
and write the comoving sound horizon as an integral over scale factor.
\end{exercise}

\begin{exercise}[Photon diffusion damping]{Cambridge Part III Cosmology, examples}
Estimate the comoving photon diffusion length before recombination as a random
walk with mean free path \((a n_e\sigma_T)^{-1}\).  Explain why perturbations
below this scale are exponentially damped.
\end{exercise}

\begin{exercise}[Spherical collapse threshold]{Tong Cosmology, Cambridge Part III Sheet 4}
Follow a uniform overdense top-hat as a closed Friedmann region in an
Einstein--de Sitter background.  Show that the linearly extrapolated density
contrast at collapse is
\(\delta_c=3(12\pi)^{2/3}/20\simeq1.686\).
\end{exercise}

\begin{exercise}[Slow-roll e-folds]{Cambridge Part III Field Theory in Cosmology, examples}
For a canonical inflaton with potential \(V(\phi)\), derive the slow-roll
parameters and the number of e-folds.  Apply the result to
\(V=\tfrac12m^2\phi^2\).
\end{exercise}

\begin{exercise}[Scalar fluctuation amplitude]{Cambridge Part III Field Theory in Cosmology, 2021 exam}
Combine the vacuum fluctuation of a light scalar at horizon exit with the
background slow-roll motion to derive the curvature power spectrum and its
spectral tilt to leading order.
\end{exercise}

\begin{exercise}[Tensor spectrum and the consistency relation]{Cambridge Part III Field Theory in Cosmology, examples}
Derive the primordial tensor power and tensor tilt in slow-roll inflation.
Combine them with the scalar result to obtain the single-field consistency
relation.
\end{exercise}
